\section{Results}

\label{sec:results}
% ══════════════════════════════════════════════════════════════════════════════

\subsection{Current Distribution Accuracy}

\begin{table}[H]
\centering
\caption{Model performance on current species distributions (mean
$\pm$ SD across 847 species). Best results in \textbf{bold}.}
\label{tab:current}
\begin{tabular}{lccc}
\toprule
\textbf{Model} & \textbf{AUC} & \textbf{TSS} & \textbf{Boyce} \\
\midrule
MaxEnt          & 0.867 $\pm$ 0.071 & 0.612 $\pm$ 0.114 & 0.751 $\pm$ 0.098 \\
Random Forest   & 0.881 $\pm$ 0.063 & 0.641 $\pm$ 0.107 & 0.773 $\pm$ 0.091 \\
BRT             & 0.889 $\pm$ 0.058 & 0.658 $\pm$ 0.102 & 0.789 $\pm$ 0.087 \\
DeepSDM         & 0.901 $\pm$ 0.054 & 0.687 $\pm$ 0.096 & 0.812 $\pm$ 0.079 \\
CNN-SDM         & 0.912 $\pm$ 0.049 & 0.711 $\pm$ 0.091 & 0.834 $\pm$ 0.072 \\
\textbf{HabitatNet} & \textbf{0.943 $\pm$ 0.038} & \textbf{0.762 $\pm$ 0.078} & \textbf{0.891 $\pm$ 0.056} \\
\bottomrule
\end{tabular}
\end{table}

HabitatNet achieves a mean AUC of 0.943, representing a 3.4\% absolute
improvement over the next-best model (CNN-SDM). The improvement is most
pronounced for species with fragmented ranges, where graph-based spatial
reasoning provides the greatest advantage.

\subsection{Temporal Transferability}

\begin{table}[H]
\centering
\caption{Temporal validation: models trained on 2000--2015, evaluated on
2016--2023 resurvey data.}
\label{tab:temporal}
\begin{tabular}{lccc}
\toprule
\textbf{Model} & \textbf{AUC} & \textbf{TSS} & \textbf{$\Delta$ AUC} \\
\midrule
MaxEnt          & 0.793 $\pm$ 0.089 & 0.497 $\pm$ 0.131 & $-$0.074 \\
Random Forest   & 0.801 $\pm$ 0.082 & 0.513 $\pm$ 0.124 & $-$0.080 \\
BRT             & 0.811 $\pm$ 0.077 & 0.531 $\pm$ 0.119 & $-$0.078 \\
DeepSDM         & 0.834 $\pm$ 0.069 & 0.567 $\pm$ 0.108 & $-$0.067 \\
CNN-SDM         & 0.849 $\pm$ 0.063 & 0.591 $\pm$ 0.101 & $-$0.063 \\
\textbf{HabitatNet} & \textbf{0.907 $\pm$ 0.047} & \textbf{0.689 $\pm$ 0.084} & $-$\textbf{0.036} \\
\bottomrule
\end{tabular}
\end{table}

HabitatNet exhibits the smallest performance degradation under temporal
transfer ($\Delta$AUC = $-$0.036), suggesting superior generalization to
novel environmental conditions. The graph structure allows the model to
propagate information about distributional shifts through connected
landscape elements.

\subsection{Climate Projection Results}

Under SSP2-4.5 (intermediate emissions), HabitatNet projects:
\begin{itemize}
  \item \textbf{Range contractions}: 67\% of species lose $>$10\% of
    current suitable habitat by 2081--2100.
  \item \textbf{Range shifts}: Mean poleward shift of 6.1\,km/decade
    ($\pm$2.3), consistent with observed trends
    \citep{chen2011rapid}.
  \item \textbf{Elevational shifts}: Mean upslope shift of 11.2\,m/decade
    in mountainous regions.
  \item \textbf{Novel disconnections}: 23\% of currently connected
    habitat networks become fragmented by 2061--2080.
\end{itemize}

Under SSP5-8.5 (high emissions), projections intensify dramatically:
range contractions affect 89\% of species, and 41\% of habitat networks
become disconnected.

\subsection{Connectivity Analysis}

\begin{table}[H]
\centering
\caption{Habitat connectivity metrics under current and projected
conditions (SSP2-4.5, 2081--2100). Mean across species.}
\label{tab:connectivity}
\begin{tabular}{lcc}
\toprule
\textbf{Metric} & \textbf{Current} & \textbf{2081--2100} \\
\midrule
Number of patches         & 4.2 $\pm$ 2.1  & 7.8 $\pm$ 3.6 \\
Largest patch area (km$^2$) & 84{,}300 $\pm$ 41{,}200 & 52{,}100 $\pm$ 29{,}800 \\
Effective mesh size       & 0.71 $\pm$ 0.18 & 0.43 $\pm$ 0.21 \\
Graph connectivity index  & 0.84 $\pm$ 0.11 & 0.58 $\pm$ 0.19 \\
\bottomrule
\end{tabular}
\end{table}

\subsection{Corridor Identification}

The greedy corridor algorithm (\Cref{alg:corridor}) identifies
restoration priorities that maximize connectivity recovery. With a
budget of 5\% additional restored area, corridors recover 72\% of lost
connectivity on average, compared to 31\% for random placement and 54\%
for least-cost path heuristics.


% ══════════════════════════════════════════════════════════════════════════════
